\documentclass[14pt,dvipdfmx]{jsarticle}
\input{lecture-preamble.tex}

% 割線 AP（平均変化率）
\newcommand{\secantfig}{%
  \begin{tikzpicture}[x=1.05cm, y=1.18cm, line cap=round]
    \draw[-{Stealth}] (-0.35,0) -- (3.35,0) node[right] {$x$};
    \draw[-{Stealth}] (0,-0.28) -- (0,2.85) node[above] {$y$};
    \node[below left, inner sep=1pt, font=\footnotesize] at (0,0) {O};
    \draw[thick, domain=0.15:2.95, samples=40] plot (\x, {0.30*\x*\x});
    \node[font=\footnotesize, anchor=south west] at (1.85,2.65) {$y=f(x)$};
    \coordinate (A) at (1.05, {0.30*1.05*1.05});
    \coordinate (P) at (2.35, {0.30*2.35*2.35});
    \draw ($(A)!-0.22!(P)$) -- ($(P)!-0.35!(A)$);
    \fill (A) circle (1.8pt);
    \fill (P) circle (1.8pt);
    \node[above=2pt, font=\footnotesize] at (A) {A};
    \node[above=2pt, font=\footnotesize] at (P) {P};
    \draw[dashed, gray] (A) -- (1.05,0) node[below, font=\footnotesize] {$a$};
    \draw[dashed, gray] (P) -- (2.35,0) node[below, font=\footnotesize] {$a+h$};
    \draw[dashed, gray] (A) -- (A -| P);
    \node[right, font=\footnotesize, inner sep=1pt]
      at ($(P)!0.5!(A -| P)$) {$f(a+h)-f(a)$};
    \node[above, font=\footnotesize, inner sep=1pt] at (1.70,0) {$h$};
  \end{tikzpicture}%
}

% 接線（h→0 で割線が接線に近づく）
\newcommand{\tangentfig}{%
  \begin{tikzpicture}[x=1.05cm, y=1.18cm, line cap=round]
    \draw[-{Stealth}] (-0.35,0) -- (3.35,0) node[right] {$x$};
    \draw[-{Stealth}] (0,-0.28) -- (0,2.85) node[above] {$y$};
    \node[below left, inner sep=1pt, font=\footnotesize] at (0,0) {O};
    \draw[thick, domain=0.15:2.95, samples=40] plot (\x, {0.30*\x*\x});
    \node[font=\footnotesize, anchor=south west] at (1.95,2.55) {$y=f(x)$};
    \coordinate (A) at (1.05, {0.30*1.05*1.05});
    % y=0.30x^2 の x=1.05 における傾きは 0.63
    \coordinate (P) at (2.45, {0.30*2.45*2.45});
    \coordinate (Q) at (1.75, {0.30*1.75*1.75});
    \draw[gray] ($(A)!-0.12!(Q)$) -- (Q);
    \draw[gray] ($(A)!-0.08!(P)$) -- (P);
    \fill (Q) circle (1.2pt);
    \fill (P) circle (1.2pt);
    \node[above=2pt, font=\footnotesize] at (P) {P};
    \draw[thick] ($(A)+(-0.42,-0.265)$) -- ($(A)+(1.85,1.166)$);
    \node[font=\footnotesize, above right] at ($(A)+(2,1.25)$) {$\ell$};
    \fill (A) circle (1.8pt);
    \node[above=2pt, font=\footnotesize] at (A) {A};
    \draw[dashed, gray] (A) -- (1.05,0) node[below, font=\footnotesize] {$a$};
    % 底辺 1、高さは傾き 0.63
    \coordinate (S) at ($(A)+(0.28,0.176)$);
    \coordinate (R) at ($(S)+(1,0)$);
    \coordinate (U) at ($(R)+(0,0.63)$);
    \draw[dashed] (S) -- (R) node[midway, below, font=\footnotesize] {$1$};
    \draw[dashed] (R) -- (U) node[midway, right, font=\footnotesize] {$f'(a)$};
  \end{tikzpicture}%
}

\begin{document}

\lectureheader{数学Ⅱ}{6}{微分法と積分法}{1}{微分係数と導関数}{1-CD}{微分係数, 図形的な意味}
  {180,181}{微分係数を求め, 接線の傾きとの関係を理解しよう}

% 1-1-AB の定義 6.1.1、6.1.2 の続き
\setcounter{definition}{2}

% ============ 1ページ目（表・左） ============
\vspace{-2.0zw}%
\heading{〜前時の復習〜}
\vspace{-0.35em}
\begin{reviewbox}{平均変化率と極限値}
  \begin{enumerate}[label=\textbf{\arabic*}, leftmargin=1.5zw, itemsep=0.2em]
    \item 平均変化率は {\renewcommand{\fitblankpad}{6zw}%
          \fitblankbf{\textstyle\frac{f(b)-f(a)}{b-a}}} で,
          \par\vspace{0.5em}
          直線の\fitblankbf{\text{傾き}}を表す.
    \item $2 + h$ の $h$ が $0$ に限りなく近づくとき, 極限値は $2$ である.
          \par\vspace{0.5em}
          これを {\renewcommand{\fitblankpad}{5zw}%
          \fitblankbf{\lim_{h \to 0}(2+h)=2}} と書く.
  \end{enumerate}
\end{reviewbox}

\vspace{0.5em}
\heading{微分係数}
\vspace{-0.3em}
\begin{definitionbox}{微分係数}
  関数 $f(x)$ の, $x=a$ から $x=a+h$ までの平均変化率
  \[
    \frac{f(a+h)-f(a)}{h}
  \]
  において, $h$ が $0$ に限りなく近づくとき,\\
  この平均変化率が一定の値に限りなく近づくならば, その極限値を,\\
  関数 $f(x)$ の $x=a$ における\fitblankbf{\text{微分係数}}, または変化率
  といい, \fitblankbf{f'(a)} で表す.
  \[
    f'(a)=\fitblankbf{\lim_{h \to 0}\frac{f(a+h)-f(a)}{h}}
  \]
\end{definitionbox}

\vfill
\nextpage

% ============ 2ページ目（表・右） ============
\vspace{-1.4zw}%
\noindent{\bfseries 例\,3}\quad
関数 $f(x)=x^{2}$ の、$x=3$ における微分係数を求める。

\vspace{0.15em}
\begin{answerbox}{}[15]
  $f'(3)=\lim_{h \to 0}\dfrac{f(3+h)-f(3)}{h}
    =\lim_{h \to 0}\dfrac{(3+h)^{2}-3^{2}}{h}$\\[0.55em]
  $=\lim_{h \to 0}\dfrac{6h+h^{2}}{h}
    =\lim_{h \to 0}\dfrac{h(6+h)}{h}$\\[0.55em]
  $=\lim_{h \to 0}(6+h)=6$
  \boxanswer{$6$}
\end{answerbox}

\begin{practiceprob}{3}
  次の微分係数を求めよ。
  \begin{tasks}[label=(\arabic*), label-width=1.6zw, item-indent=2zw, after-item-skip=0.15em](1)
    \task 関数 $f(x)=x^{2}$ の、$x=1$ における微分係数
    \task 関数 $f(x)=3x^{2}$ の、$x=-2$ における微分係数
  \end{tasks}
\end{practiceprob}
\vspace{-0.15em}
\begin{answerbox}{(1)}[15]
  $f'(1)=\lim_{h \to 0}\dfrac{(1+h)^{2}-1^{2}}{h}
    =\lim_{h \to 0}\dfrac{2h+h^{2}}{h}$\\[0.55em]
  $=\lim_{h \to 0}(2+h)=2$
  \boxanswer{$2$}
\end{answerbox}
\begin{answerbox}{(2)}[15]
  $f'(-2)=\lim_{h \to 0}\dfrac{3(-2+h)^{2}-3\cdot(-2)^{2}}{h}$\\[0.55em]
  $=\lim_{h \to 0}\dfrac{3(4-4h+h^{2})-12}{h}
    =\lim_{h \to 0}\dfrac{-12h+3h^{2}}{h}$\\[0.55em]
  $=\lim_{h \to 0}(-12+3h)=-12$
  \boxanswer{$-12$}
\end{answerbox}

\vfill
\nextpage

% ============ 3ページ目（裏・左） ============
\vspace{-1.4zw}%
\heading{微分係数の図形的な意味}
\vspace{-0.2em}
関数 $f(x)$ が微分係数 $f'(a)$ をもつとする.
\par\vspace{0.35em}
\sidebyside[0.56]{%
  関数 $y=f(x)$ のグラフ上に $2$ 点
  $\mathrm{A}(a,\ f(a))$, $\mathrm{P}(a+h,\ f(a+h))$
  をとると, 直線 $\mathrm{AP}$ の傾き
  \[
    \frac{f(a+h)-f(a)}{h}
  \]
  は, $x=a$ から $x=a+h$ までの平均変化率に等しい.
}{\secantfig}

\vspace{0.45em}
\sidebyside[0.56]{%
  $h$ が $0$ に限りなく近づくとき, 点 $\mathrm{P}$ は点 $\mathrm{A}$ に
  限りなく近づく. このとき,
  \[
    \lim_{h \to 0}\frac{f(a+h)-f(a)}{h}=f'(a)
  \]
  であるから, 直線 $\mathrm{AP}$ は点 $\mathrm{A}$ を通り, 傾きが
  $f'(a)$ の直線 $\ell$ に限りなく近づく.
}{\tangentfig}

\vspace{0.45em}
この直線 $\ell$ を, グラフ上の点 $\mathrm{A}$ における\fitblankbf{\text{接線}}といい,\\
$\mathrm{A}$ を\fitblankbf{\text{接点}}という.\\
また, 直線 $\ell$ はこの曲線に点 $\mathrm{A}$ で\fitblankbf{\text{接する}}という.

\vspace{0.5em}
\begin{definitionbox}{接線の傾き}
  関数 $y=f(x)$ のグラフ上の点 $\mathrm{A}(a,\ f(a))$ における接線の傾きは,
  関数 $f(x)$ の $x=a$ における微分係数 \fitblankbf{f'(a)} に等しい.
\end{definitionbox}

\vfill
\nextpage

% ============ 4ページ目（裏・右） ============
\vspace{-1.4zw}%
\noindent{\bfseries 例\,4}\quad
関数 $y=x^{2}$ のグラフ上の点 $(3,\ 9)$ における接線の傾き $m$ を求める。

\vspace{0.15em}
\begin{answerbox}{}[10]
  $f(x)=x^{2}$ とすると, $m=f'(3)$.\\[0.4em]
  例3より, $f'(3)=6$ であるから, $m=6$.
  \boxanswer{$6$}
\end{answerbox}

\begin{practiceprob}{4}
  関数 $y=x^{2}$ のグラフ上の次の点における接線の傾きを求めよ。
  \begin{tasks}[label=(\arabic*), label-width=1.6zw, item-indent=2zw, after-item-skip=0.15em](1)
    \task 点 $(1,\ 1)$
    \task 点 $(-2,\ 4)$
  \end{tasks}
\end{practiceprob}
\vspace{-0.15em}
\begin{answerbox}{(1)}[8]
  練習3\,(1) より, $f'(1)=2$ であるから, 傾きは $2$.
  \boxanswer{$2$}
\end{answerbox}
\begin{answerbox}{(2)}[14]
  $f'(-2)=\lim_{h \to 0}\dfrac{(-2+h)^{2}-(-2)^{2}}{h}$\\[0.5em]
  $=\lim_{h \to 0}\dfrac{-4h+h^{2}}{h}=\lim_{h \to 0}(-4+h)=-4$
  \boxanswer{$-4$}
\end{answerbox}

\vfill
\begin{summarybox}
  \begin{enumerate}[label=\textbf{\arabic*}, leftmargin=1.5zw, itemsep=0.35em]
    \item $x=a$ における微分係数は
          \par\vspace{0.35em}
          {\renewcommand{\fitblankpad}{3zw}%
          \fitblankbf{f'(a)=\lim_{h \to 0}\dfrac{f(a+h)-f(a)}{h}}}
          \mbox{である.}
    \item 接線の傾きは, 微分係数 \fitblankbf{f'(a)} に等しい.
  \end{enumerate}
\end{summarybox}

\vspace{0.4em}
\nexttime{各点の微分係数を対応させて, 導関数を考えてみよう!}

\end{document}
